\chapter{Machine-Learning and Data Engineer}\label{ch:in:machine-learning-and-data-engineer}

In fiscal 2021 the sourced finance employers filed 136 labour condition applications for jobs titled in machine learning,
data science or data engineering; in fiscal 2025 they filed 1\,051, almost eight times as many. Over the same four years
their filings for software engineers doubled and those for \omterm{def:in:quant-researcher:def}{quantitative researchers} did not move. The occupation that most
of the 2025 filings name, \omterm{def:in:machine-learning-and-data-engineer:roles}{data scientist}, did not exist in the US classification of 2010; it was created in the 2018
revision. This chapter describes the machine-learning and data jobs in trading and finance, measures how fast their filings
grew and with what uncertainty, and compares their pay with the software engineers' beside them.

\begin{center}\footnotesize
\begin{tabular}{@{}p{2.4cm}p{3.3cm}p{3.3cm}p{3.3cm}@{}}
\toprule
\multicolumn{4}{@{}l}{\textbf{Role cards: machine learning and data}}\\
\midrule
 & \omterm{def:in:machine-learning-and-data-engineer:roles}{data scientist} & machine-learning engineer & \omterm{def:in:machine-learning-and-data-engineer:roles}{data engineer}\\
\midrule
works on & models from data, for a research or business question & the systems that train, serve and monitor models & the pipelines and stores that data lives in\\
sits in & research groups or business lines & a platform team or a research group & a data or platform team\\
codes in filings & 15-2051 & 15-2051, 15-1252 & 15-1252, 15-1243\\
taught in & Book~12, ch.~1; Book~7, ch.~12 & Book~12, ch.~23--28 & Book~15, ch.~2--4; Book~12, ch.~24\\
\bottomrule
\end{tabular}
\end{center}

\begin{definition}[Data scientist, data engineer]\label{def:in:machine-learning-and-data-engineer:roles}
A \emph{data scientist}\index{data scientist} is an analyst who builds statistical and machine-learning models from data
to answer a question or to drive a decision, and reports what the data show. A \emph{data engineer}\index{data engineer}
is an engineer who builds and runs the pipelines, storage and interfaces through which data is acquired, cleaned,
versioned and served to the people and systems that use it.
\end{definition}

The occupational definition of the first is broad: \omterm{def:in:machine-learning-and-data-engineer:roles}{data scientists} ``Develop and implement a set of techniques or
analytics applications to transform raw data into meaningful information'' and ``Apply data mining, data modeling,
natural language processing, and machine learning''. The classification has no occupation for the second: O*NET lists
``Data Engineer'' among the titles reported for database architects (15-1243), and employers file most \omterm{def:in:machine-learning-and-data-engineer:roles}{data engineers} as
software developers. The machine-learning engineer, who builds the systems that train, serve and monitor models, is
Book~12's (chapter~28).

\section{Research-embedded machine-learning roles}

In a trading firm's research group, machine learning is a method, not a department. The researcher of chapter~17 who fits
a gradient-boosted model to order-book features is doing machine learning; the title on the contract may say
\omterm{def:in:quant-researcher:def}{quantitative researcher} or \omterm{def:in:machine-learning-and-data-engineer:roles}{data scientist}, and the filings follow the title (chapter~19). What distinguishes the
research-embedded role is the problem: a signal with a low signal-to-noise ratio, few independent observations,
non-stationary data and a cost of being wrong that is paid in money. Book~12 is about exactly that difference
(chapter~1), and its tools (validation that respects time, labels built from future prices, models that survive noise)
are the role's daily work.

A research-embedded machine-learning specialist is judged as a researcher is: on out-of-sample results, and on whether the
model adds to what simpler models already do. Book~12's habit of setting every model against a linear and regularised
baseline (chapter~4) is the role's discipline. The \omterm{def:in:quant-researcher:feedback}{feedback speed} of chapter~17 applies unchanged: a model rebalanced monthly
earns evidence as slowly as any monthly signal.

\section{Platform machine-learning roles}

A firm that runs many models needs a machine-learning platform (Book~12, chapter~28): data and feature stores, training
infrastructure, experiment tracking and a model registry, compilation and serving, monitoring and retraining. The
machine-learning engineer builds and runs it, and the handoff specification of Book~12 is the contract between the
people who build models and the people who run them. Most of the growth in the filings (below) is at banks, and the
filings do not say which of a bank's businesses each job serves.

The platform role is judged like the \omterm{def:in:quant-developer:research}{research engineer}'s of chapter~19: on how fast and how safely models reach
production and on how rarely they fail there. It is closer to software engineering than to research, and its pay and
titles follow software engineering's more than research's.

\section{The data engineer}

A trading firm is a consumer of data before it is anything else: market data by venue and by day, reference data,
corporate actions, alternative data (Book~7, chapter~12), its own orders and fills. The \omterm{def:in:machine-learning-and-data-engineer:roles}{data engineer} owns the path from
the vendor's or venue's file to the researcher's query: capture and storage of ticks (Book~15, chapter~2), columnar
formats and tick stores (chapters~3 and~4), point-in-time versions that let a backtest see only what was known at the
time, feature stores that serve the same values to training and to production (Book~12, chapter~24), and the checks that
catch a vendor's silent change before a model trains on it.

The work is unglamorous and indispensable: a research result is only as good as the data under it, and a missing day, a
wrong timestamp or a survivor-biased universe can produce a backtest that is false in every detail. \omterm{def:in:machine-learning-and-data-engineer:roles}{Data engineers} are
judged on completeness, correctness and timeliness of what they serve, and on the cost of storing and moving it.

\section{Growth and pay}

\begin{dated}{2025-09}{Machine-learning and data roles in the filings}\label{dat:in:machine-learning-and-data-engineer:filings}
Labour condition applications at the sourced finance employers, fiscal 2021 and 2025: machine-learning and data titles
136 and 1\,051 (banks 920 in 2025); software engineers 2\,618 and 5\,727; \omterm{def:in:quant-researcher:def}{quantitative researchers} 1\,287 and 1\,297;
all families 4\,422 and 8\,808. Fiscal 2025, \omterm{def:in:machine-learning-and-data-engineer:roles}{data scientists} (15-2051) against software developers (15-1252) at the same
employers: median offered base \$138\,000 against \$155\,000; by \omterm{def:in:pay-levels-by-role-firm-type-and-seniority:soc}{wage level} $-\$20\,500$ (I), $-\$18\,300$ (II),
$-\$19\,000$ (III), $+\$2\,220$ (IV). \omterm{def:in:how-pay-works:base}{Base salary} only.
\end{dated}

\begin{method}[Growth of counts]\label{met:in:machine-learning-and-data-engineer:trend}
Let $n_t$ be the filings in year $t$. Model them as Poisson with mean $\lambda_0e^{g t}$; the weighted least-squares fit
of $\log n_t$ on $t$ with weights $n_t$ estimates the continuous growth rate $g$ with standard error
$1/\sqrt{\sum_t n_t(t-\bar t)^2}$. With two years $t_0,t_1$ it is $g=\log(n_1/n_0)/(t_1-t_0)$ with standard error
$\sqrt{1/n_0+1/n_1}/(t_1-t_0)$, and the annual rate is $e^g-1$. Each count has an exact Poisson interval.
\end{method}

The machine-learning and data filings grew at 66.7\% a year (95\% interval 59.4\% to 74.3\%), the software engineers' at
21.6\% (20.2\% to 23.0\%); the \omterm{def:in:quant-researcher:def}{quantitative researchers}' did not grow (0.2\%, $-1.7\%$ to 2.1\%) and the traders' fell
slightly (\cref{fig:in:machine-learning-and-data-engineer:growth}). The family's share of all filings rose from 3.1\% to
11.9\%. The Poisson interval is the smallest honest one: filings come in batches from a few employers, and nine in ten of
the 2025 machine-learning filings came from banks, so the true uncertainty about the industry's hiring is larger than the
interval says.

\begin{omfigure}
\begin{tikzpicture}
\begin{axis}[ybar,area legend,width=10.5cm,height=5.4cm,xmin=0.5,xmax=6.5,ymode=log,ymin=10,ymax=10000,log origin=infty,bar width=11pt,
  xtick={1,2,3,4,5,6},xticklabels from table={figdata/industry/21-machine-learning-and-data-engineer/counts.csv}{family},
  xticklabel style={font=\scriptsize,text width=1.6cm,align=center},ylabel={\small applications (log scale)},
  yticklabel style={font=\footnotesize},grid=major,grid style={gray!20},table/col sep=comma,
  legend style={legend cell align=left,font=\scriptsize,at={(0.5,-0.42)},anchor=north,legend columns=2,/tikz/every even column/.append style={column sep=8pt}}]
\addplot[fill=black!30,draw=none,error bars/.cd,y dir=both,y explicit]
  table[x=pos,y=n2021,y error plus expr=\thisrow{hi2021}-\thisrow{n2021},y error minus expr=\thisrow{n2021}-\thisrow{lo2021}]
  {figdata/industry/21-machine-learning-and-data-engineer/counts.csv};
\addplot[fill=omDef,draw=none,error bars/.cd,y dir=both,y explicit]
  table[x=pos,y=n2025,y error plus expr=\thisrow{hi2025}-\thisrow{n2025},y error minus expr=\thisrow{n2025}-\thisrow{lo2025}]
  {figdata/industry/21-machine-learning-and-data-engineer/counts.csv};
\legend{fiscal 2021,fiscal 2025}
\end{axis}
\end{tikzpicture}

\omcaption{Labour condition applications by role family at the sourced finance employers, fiscal 2021 and 2025, with
exact 95\% Poisson intervals (barely visible at these counts). Data: \texttt{data/industry/lca\_ranges.csv}, through
\texttt{in\_mldata.counts}.}\label{fig:in:machine-learning-and-data-engineer:counts}
\end{omfigure}

\begin{omfigure}
\begin{tikzpicture}
\begin{axis}[width=10.5cm,height=5cm,xmin=-20,xmax=90,ymin=0.4,ymax=6.6,ytick=data,
  yticklabels from table={figdata/industry/21-machine-learning-and-data-engineer/growth.csv}{family},
  yticklabel style={font=\footnotesize},xticklabel style={font=\footnotesize},
  xlabel={\small annual growth of filings, fiscal 2021 to 2025, \%},grid=major,grid style={gray!20},table/col sep=comma]
\addplot[only marks,mark=*,omDef,error bars/.cd,x dir=both,x explicit]
  table[x=annual,y=pos,x error plus expr=\thisrow{hi}-\thisrow{annual},x error minus expr=\thisrow{annual}-\thisrow{lo}]
  {figdata/industry/21-machine-learning-and-data-engineer/growth.csv};
\draw[black!50] (axis cs:0,0.4) -- (axis cs:0,6.6);
\end{axis}
\end{tikzpicture}

\omcaption{Annual growth of filings by role family between fiscal 2021 and 2025, with 95\% intervals under the Poisson
model. Data: \texttt{in\_mldata.trends}, from \texttt{firm.roles.filing\_trend}.}\label{fig:in:machine-learning-and-data-engineer:growth}
\end{omfigure}

The pay comparison is less flattering than the growth. At the same employers, filings coded as \omterm{def:in:machine-learning-and-data-engineer:roles}{data scientists} carry a
median offered base \$17\,000 below the software developers' (95\% interval $-\$20\,000$ to $-\$15\,000$), and about
\$19\,000 below at each of levels I to III; only at level IV are the two level (\cref{fig:in:machine-learning-and-data-engineer:gap}).
Most of the growth is at banks, whose filings do not say which business a job serves; the few
machine-learning filings at trading firms carry higher medians (\$175\,000 at market makers, \$193\,000 at platforms and
\$192\,912 at systematic funds, each from ten to twelve filings), in line with those firms' researchers and engineers.

\begin{omfigure}
\begin{tikzpicture}
\begin{axis}[width=10.5cm,height=5cm,xmin=0.5,xmax=5.5,ymin=-30,ymax=10,xtick={1,2,3,4,5},
  xticklabels={all levels,I,II,III,IV},xlabel={\small wage level},ylabel={\small gap, \$ thousand},
  tick label style={font=\footnotesize},grid=major,grid style={gray!20},table/col sep=comma]
\addplot[only marks,mark=*,omDef,thick,error bars/.cd,y dir=both,y explicit]
  table[x=pos,y=diff,y error plus expr=\thisrow{hi}-\thisrow{diff},y error minus expr=\thisrow{diff}-\thisrow{lo}]
  {figdata/industry/21-machine-learning-and-data-engineer/gap.csv};
\draw[black!50] (axis cs:0.5,0) -- (axis cs:5.5,0);
\end{axis}
\end{tikzpicture}

\omcaption{Median offered base of data-scientist filings minus software developers' at the sourced finance employers,
fiscal 2025, overall and by \omterm{def:in:pay-levels-by-role-firm-type-and-seniority:soc}{wage level}, with 95\% bootstrap intervals. Data: \texttt{data/industry/lca\_ds\_gap.csv},
through \texttt{in\_mldata.gaps}.}\label{fig:in:machine-learning-and-data-engineer:gap}
\end{omfigure}

\begin{dated}{2025-05}{Where data scientists work: the survey}\label{dat:in:machine-learning-and-data-engineer:survey}
Occupational survey, May 2025, \omterm{def:in:machine-learning-and-data-engineer:roles}{data scientists} (15-2051): finance and insurance 46\,730 employed, median \$124\,770;
information 32\,410, median \$141\,440; professional, scientific and technical services 69\,730, median \$126\,730.
Within finance: credit intermediation 12\,450 (median \$130\,300), insurance carriers 15\,090 (\$107\,680), securities
7\,210 (\$134\,510). Software publishers 10\,950 (\$156\,220).
\end{dated}

Finance and insurance employs more \omterm{def:in:machine-learning-and-data-engineer:roles}{data scientists} than the information sector, though at a lower median, and the
securities industry is a small part of it (\cref{fig:in:machine-learning-and-data-engineer:survey}). The picture matches
the filings: data science in finance is mostly a bank and insurance job, and in trading it is mostly one of the tools of
research.

\begin{omfigure}
\begin{tikzpicture}
\begin{axis}[xbar,width=10.5cm,height=5.4cm,xmin=0,xmax=85,ymin=0.4,ymax=7.6,bar width=9pt,ytick=data,
  yticklabels from table={figdata/industry/21-machine-learning-and-data-engineer/survey.csv}{industry},
  yticklabel style={font=\scriptsize},xticklabel style={font=\footnotesize},
  xlabel={\small data scientists employed, thousands},grid=major,grid style={gray!20},table/col sep=comma,
  nodes near coords,nodes near coords style={font=\scriptsize,anchor=west,xshift=1pt,yshift=0pt},
  point meta=explicit symbolic]
\addplot[fill=omDef,draw=none] table[x=employment,y=pos,meta=median]
  {figdata/industry/21-machine-learning-and-data-engineer/survey.csv};
\end{axis}
\end{tikzpicture}

\omcaption{\omterm{def:in:machine-learning-and-data-engineer:roles}{Data scientists} employed by sector and industry, May 2025 occupational survey; labels give the median annual
wage in \$ thousand. The top three bars are sectors; the four below are industries within them (banks, insurance
carriers and securities within finance and insurance; software publishers within information). Data:
\texttt{data/industry/oews\_roles.csv}, through \texttt{in\_mldata.survey}.}\label{fig:in:machine-learning-and-data-engineer:survey}
\end{omfigure}

\section{Tutorial: the fastest-growing title}\label{tut:in:machine-learning-and-data-engineer:tut}

\begin{tutorial}
\textbf{Goal.} Measure the growth of the machine-learning and data filings with its uncertainty, and put their pay beside
the software engineers'.
\textbf{End state:} \cref{fig:in:machine-learning-and-data-engineer:counts,fig:in:machine-learning-and-data-engineer:growth,fig:in:machine-learning-and-data-engineer:gap}.
\begin{tutsteps}
\item \textbf{Counts.} Sum chapter~14's \texttt{lca\_ranges.csv} over employer kinds for each role family and year;
\texttt{firm.roles.poisson\_interval(n)} gives each count's exact interval.
\item \textbf{Growth.} \texttt{filing\_trend(years, counts)} fits the log-linear Poisson model
(\cref{lst:in:machine-learning-and-data-engineer:trend}).
\omcode{code/firm/roles/firm_roles.py}{320}{331}{Growth of filing counts with its standard error.}\label{lst:in:machine-learning-and-data-engineer:trend}
\item \textbf{Pay.} \texttt{in\_lca\_ds\_derive.py} compares data-scientist and software-developer filings level by level
from chapter~14's cache with \texttt{median\_gap} (chapter~19).
\item \textbf{Survey.} \texttt{oews\_roles.csv} for 15-2051 by sector and industry.
\end{tutsteps}
\end{tutorial}

For the machine-learning family: $g=\log(1\,051/136)/4=0.511$ with standard error 0.023; 66.7\% a year. Its 2025 count's
Poisson interval is 988 to 1\,117.

\textbf{What to change next.} Read the other fiscal years and fit the trend over all of them; separate banks from trading
firms; replace the Poisson model by one that allows for employers filing in batches (a negative binomial, or a bootstrap
over employers).

\section{Build: machine-learning cards and filing trends}\label{bld:in:machine-learning-and-data-engineer:roles}

\begin{build}
\bfield{Purpose} Add the machine-learning and data roles to the registry, and measure the growth of a count with its
uncertainty.
\bfield{Interface} \texttt{firm.roles}: the cards \texttt{data scientist}, \texttt{data engineer},
\texttt{machine-learning engineer}; \texttt{poisson\_interval(n, level) -> (lo, hi)};
\texttt{filing\_trend(years, counts, z) -> dict(rate, se, annual, lo, hi)}.
\bfield{Rules} Counts are Poisson unless the caller says otherwise; the interval is exact (Garwood); growth is
continuous in the fit and reported as an annual compound rate.
\bfield{Acceptance tests} \texttt{code/firm/roles/tests/}: two points give the closed form; counts doubling each year give
$\log 2$ exactly; the interval for ten is 4.80 to 18.39; zero gives a lower bound of zero.
\bfield{Stretch} Overdispersion; a trend per employer kind with a test of equal growth; forecasts with intervals.
\end{build}

\begin{omsources}
\item Bureau of Labor Statistics, SOC 2010 to 2018 crosswalk; occupational survey, May 2025.
\item O*NET OnLine, occupations 15-2051 and 15-1243 (CC BY 4.0).
\item Chapter~14's tables from the Department of Labor's LCA files.
\item Garwood, F. (1936), Fiducial limits for the Poisson distribution, \emph{Biometrika} 28(3/4), 437--442.
\end{omsources}

\section{Exercises}

\begin{exercise}[$\star$]\label{exo:in:machine-learning-and-data-engineer:1}
By what factor did the machine-learning and data filings grow between fiscal 2021 and 2025, and what is that as an annual
rate?
\end{exercise}

\begin{exercise}[$\star$]\label{exo:in:machine-learning-and-data-engineer:2}
What share of all filings did the family have in each year?
\end{exercise}

\begin{exercise}[$\star$]\label{exo:in:machine-learning-and-data-engineer:3}
From the survey box, which employs more \omterm{def:in:machine-learning-and-data-engineer:roles}{data scientists}, finance and insurance or the information sector, and which pays
the higher median?
\end{exercise}

\begin{exercise}[$\star\star$]\label{exo:in:machine-learning-and-data-engineer:4}
Compute the standard error of the growth rate for the software engineers and explain why it is smaller than the
machine-learning family's.
\end{exercise}

\begin{exercise}[$\star\star$]\label{exo:in:machine-learning-and-data-engineer:5}
Why can the growth of filings coded 15-2051 not be measured from fiscal 2021, although the growth of the title family can?
\end{exercise}

\begin{exercise}[$\star\star$]\label{exo:in:machine-learning-and-data-engineer:6}
A graduate compares a data-scientist offer at a bank with a quantitative-researcher offer at a systematic fund. What does
the chapter's evidence say, and what does it not?
\end{exercise}

\begin{exercise}[$\star\star\star$]\label{exo:in:machine-learning-and-data-engineer:7}
\emph{Coding.} Suppose both years' machine-learning filings came in batches of ten, so that each count behaves like ten
times a Poisson count. Recompute the growth rate's standard error and interval.
\end{exercise}

\begin{exercise}[$\star\star\star$]\label{exo:in:machine-learning-and-data-engineer:8}
\emph{Find the flaw.} ``Machine-learning filings grew 67\% a year while researchers' did not grow, so trading firms are
replacing \omterm{def:in:quant-researcher:def}{quantitative researchers} with \omterm{def:in:machine-learning-and-data-engineer:roles}{data scientists}.''
\end{exercise}

\section{Problem: The Fastest-Growing Title}

\begin{problem}[{Weekend problem --- the fastest-growing title}]\label{pb:in:machine-learning-and-data-engineer:1}
A student deciding between a statistics master's with a machine-learning track and a financial mathematics master's asks
which way the industry's hiring is moving.

\medskip\noindent\textbf{Part I --- The roles.}
\begin{enumerate}
\item Define the \omterm{def:in:machine-learning-and-data-engineer:roles}{data scientist} and the \omterm{def:in:machine-learning-and-data-engineer:roles}{data engineer}.
\item What distinguishes a research-embedded machine-learning role?
\item What does the machine-learning engineer build, and what is the handoff specification?
\item What does a \omterm{def:in:machine-learning-and-data-engineer:roles}{data engineer} own at a trading firm?
\item How does the classification treat \omterm{def:in:machine-learning-and-data-engineer:roles}{data engineers}?
\end{enumerate}

\medskip\noindent\textbf{Part II --- The counts.}
\begin{enumerate}[resume]
\item Give the family counts for fiscal 2021 and 2025.
\item State the growth model and its standard error.
\item Give the annual growth rates with intervals for the machine-learning family, software engineers and researchers.
\item Why is the Poisson interval the smallest honest interval?
\item Which kind of employer filed most of the 2025 machine-learning applications?
\end{enumerate}

\medskip\noindent\textbf{Part III --- The pay.}
\begin{enumerate}[resume]
\item Give the data-scientist gap against software developers overall and by level.
\item What do the trading firms' few machine-learning filings show?
\item Where do \omterm{def:in:machine-learning-and-data-engineer:roles}{data scientists} work, by the survey?
\item Why might a bank's \omterm{def:in:machine-learning-and-data-engineer:roles}{data scientist} be paid less than its software developer at the same level?
\item What does \omterm{def:in:how-pay-works:base}{base salary} leave out here?
\end{enumerate}

\medskip\noindent\textbf{Part IV --- The verdict.}
\begin{enumerate}[resume]
\item State the \emph{named result}: the annual growth rate of machine-learning and data filings at finance employers,
fiscal 2021 to 2025, with its interval, and the median-base gap to software developers at the same \omterm{def:in:pay-levels-by-role-firm-type-and-seniority:soc}{wage level}.
\item Is the growth about trading?
\item Which master's prepares for which jobs?
\item What should the student ask an employer about a data-science role?
\item In two sentences, answer the student.
\end{enumerate}
\end{problem}

\section{Interview questions}

\begin{interviewq}[$\star$ \iqroles{mle, researcher}]\label{iq:in:machine-learning-and-data-engineer:1}
Why is shuffled k-fold cross-validation wrong for a daily return-prediction model?
\end{interviewq}

\begin{interviewq}[$\star$ \iqroles{mle}]\label{iq:in:machine-learning-and-data-engineer:2}
A model's live accuracy is lower than its validation accuracy. List five causes in order of how you would check them.
\end{interviewq}

\begin{interviewq}[$\star\star$ \iqroles{mle, developer}]\label{iq:in:machine-learning-and-data-engineer:3}
How do you make sure a feature has the same value in training and in production?
\end{interviewq}

\begin{interviewq}[$\star\star$ \iqroles{developer}]\label{iq:in:machine-learning-and-data-engineer:4}
A vendor silently changes the definition of a field. How would your pipeline notice?
\end{interviewq}

\begin{interviewq}[$\star\star$ \iqroles{mle, researcher}]\label{iq:in:machine-learning-and-data-engineer:5}
Your gradient-boosted model beats a linear model by 0.3\% of accuracy. Is it worth deploying?
\end{interviewq}

\begin{interviewq}[$\star\star\star$ \iqroles{mle}]\label{iq:in:machine-learning-and-data-engineer:6}
Design the storage for ten years of order-book data that researchers query by symbol and time range.
\end{interviewq}
